\chapter{シミュレーション結果：平面モデル（Sim01, Sim02）}\label{ch:planar}

モデルは、人間を矢状面（体を左右に分ける面）に投影した粗い平面モデルである（胴体 60 kg, 0.35×1.0 m / 腕 4 kg, 0.7 m / 脚 12 kg, 0.9 m）。図~\ref{fig:tikz-planar} に構成を示す。

\begin{figure}[htbp]\centering
\resizebox{0.75\linewidth}{!}{% 平面モデル：胴体＋腕＋脚
\begin{tikzpicture}[>=Stealth,scale=1.15,font=\small]
\begin{scope}[rotate=12]
  \draw[thick,fill=gray!15] (-0.35,-1) rectangle (0.35,1);
  \fill (0,0) circle (1.5pt) node[left=6pt,font=\footnotesize]{胴体重心};
  \draw[dashed,gray] (0,-1.4)--(0,1.5);
  \coordinate (S) at (0,0.7); \coordinate (H) at (0,-0.3);
  \draw[line width=3pt,blue!60] (S) -- ++(35:1.4) coordinate (A);
  \draw[line width=4pt,red!60] (H) -- ++(-60:1.8) coordinate (L);
  \fill (S) circle (2pt); \fill (H) circle (2pt);
  \draw[->] (S)++(0.7,0) arc (0:35:0.7); \node at ($(S)+(10:0.95)$) {$q_1$};
  \draw[gray] (S)--++(1,0); \draw[gray] (H)--++(1,0);
  \draw[->] (H)++(0.8,0) arc (0:-60:0.8); \node at ($(H)+(-25:1.05)$) {$q_2$};
  \node[blue!70!black,above] at (A) {腕 $m_1,\ \ell_1$};
  \node[red!70!black,right] at (L) {脚 $m_2,\ \ell_2$};
  \node[right,font=\footnotesize] at (S) {肩};
  \node[left=4pt,font=\footnotesize,yshift=-6pt] at (H) {股関節};
\end{scope}
\draw[gray] (0,0)--(0,1.8);
\draw[->] (0,1.5) arc (90:102:1.5); \node[above] at (-0.15,1.55) {$\theta$};
\node[align=left,font=\footnotesize] at (4.2,1.3) {$L = A(q)\dot\theta + B(q)\cdot\dot q$\\$L=0 \Rightarrow \dot\theta = -\Gamma(q)\cdot\dot q$};
\end{tikzpicture}
}
\caption{平面モデル：胴体＋腕＋脚。胴体の姿勢角 $\theta$、関節角 $q_1$（腕）、$q_2$（脚）。}\label{fig:tikz-planar}
\end{figure}

\section{\texorpdfstring{Sim01: 腕 1 本の往復（\texttt{sims/sim01\_single\_arm.py}）}{Sim01: 腕 1 本の往復（sims/sim01\_single\_arm.py）}}

\FIG{sim01_timeseries.png}

\begin{longtable}[]{@{}
  >{\raggedright\arraybackslash}p{(\columnwidth - 2\tabcolsep) * \real{0.5000}}
  >{\raggedright\arraybackslash}p{(\columnwidth - 2\tabcolsep) * \real{0.5000}}@{}}
\toprule\noalign{}
\begin{minipage}[b]{\linewidth}\raggedright
項目
\end{minipage} & \begin{minipage}[b]{\linewidth}\raggedright
結果
\end{minipage} \\
\midrule\noalign{}
\endhead
\bottomrule\noalign{}
\endlastfoot
腕を下→上へ 180° 振ったときの胴体回転 & \textbf{−16.0°}（逆向き） \\
腕を戻した後の残留姿勢 & 約 1e-9°（可逆的） \\
角運動量の残差 & 約 1e-5 kg·m²/s（数値微分による誤差の範囲） \\
18 m 級（長さ ×10、質量 ×1000）での同条件 & \textbf{−16.0432°（人間大と同一）} \\
\end{longtable}

\FIG{sim01_sweep.png}

\textbf{考察} - 腕 1 本では「一時的な首振り」しかできない。姿勢を保ったまま腕を戻す手段がない。 - 回転量は腕の質量比 μ と長さ比にほぼ比例して増える。μ≈0.06（人間の腕）で約 16°。 AMBAC を主機能とするなら、四肢を \textbf{重く・長く} する設計圧力がかかる（推進剤タンクを四肢に置く等）。 - スケール不変性により「巨大ロボだから AMBAC が効く/効かない」ということは \textbf{回転角については無い}。

\section{\texorpdfstring{Sim02: 腕＋脚の周期運動（\texttt{sims/sim02\_cyclic\_gait.py}）}{Sim02: 腕＋脚の周期運動（sims/sim02\_cyclic\_gait.py）}}

\FIG{sim02_cyclic.png}

\begin{longtable}[]{@{}
  >{\raggedright\arraybackslash}p{(\columnwidth - 2\tabcolsep) * \real{0.5000}}
  >{\raggedright\arraybackslash}p{(\columnwidth - 2\tabcolsep) * \real{0.5000}}@{}}
\toprule\noalign{}
\begin{minipage}[b]{\linewidth}\raggedright
項目
\end{minipage} & \begin{minipage}[b]{\linewidth}\raggedright
結果
\end{minipage} \\
\midrule\noalign{}
\endhead
\bottomrule\noalign{}
\endlastfoot
1 周あたりの正味回転（数値積分） & \textbf{+2.9778°} \\
1 周あたりの正味回転（曲率の面積分） & \textbf{+2.9778°}（一致） \\
360° 回転に必要な周回数 & 約 121 周 \\
\end{longtable}

\textbf{考察} - 形状を周期的に戻しても姿勢が累積する ＝ 非ホロノミック AMBAC が成立。 - 今回のループは曲率の弱い領域に置かれている。曲率マップ中央の負のピーク （腕 ≈ −20°、脚 ≈ +50°）を囲めば 1 周あたりの回転量は増やせる → ループ最適化が次の課題。 - 1 周あたり約 3° は人間の宇宙飛行士の経験（ゆっくりとしか向きを変えられない）と定性的に整合。
